% Part: normal-modal-logic
% Chapter: axioms-systems
% Section: derived-rules

\documentclass[../../../include/open-logic-section]{subfiles}

\begin{document}

\olfileid{nml}{prf}{der}

\olsection{القواعد المشتقة}

في طلب !!{derivation}s وكتابتها عناء وتكرار كثير.
فالانتقال من~$!A \lif !B$ إلى~$\Box !A \lif \Box !B$، وهو~\RK،
يحتاج في التفصيل إلى~\Nec، ثم إلى حالة مناسبة من~\Ax{K}،
ثم إلى~\MP. وكذلك الانتقال من~$!A \lif !B$ و$!B \lif !C$
إلى~$!A \lif !C$ يقتضي كتابة التحصيل المطوّل
\[
(!A \lif !B) \lif ((!B \lif !C) \lif (!A \lif !C))
\]
ثم إجراء~\MP{} مرتين. وقد نحتاج إلى إبدال !!{formula} جزئية
بمكافئ لها؛ كإبدال \iftag{prvDiamond}{$\Diamond
  !A$ بـ$\lnot\Box\lnot !A$}{$\Box !A$ بـ$\lnot\Diamond\lnot !A$}،
أو~$\lnot\lnot !A$ بـ$!A$. فنجعل للخطوات المتكررة أحكامًا
تختصر كتابة !!{derivation}، وتُظهر مع ذلك لماذا تكون النتائج
الوسيطة !!{derivable}. وهذا غرض الحقائق الآتية في !!{derivability}.

\begin{prop}
  إذا كان $\Log{K} \Proves !A_{\OLMathNumeral{1}}$، \dots، $\Log{K} \Proves !A_{\OLId{latin-ordinary}{n}}$،
  وكانت $!B$ تلزم عن $!A_{\OLMathNumeral{1}}$، \dots، $!A_{\OLId{latin-ordinary}{n}}$ بمنطق القضايا، فعندئذ
  $\Log{K} \Proves !B$.
\end{prop}

\begin{proof}
  إذا كانت $!B$ تلزم عن $!A_{\OLMathNumeral{1}}$، \dots،~$!A_{\OLId{latin-ordinary}{n}}$ بمنطق القضايا، فإن
  \[
  !A_{\OLMathNumeral{1}} \lif (!A_{\OLMathNumeral{2}} \lif \cdots (!A_{\OLId{latin-ordinary}{n}} \lif !B)\dots)
  \]
  مثال لقضية صادقة صدقا تحليليا. وبتطبيق \MP{} عدد ${\OLId{latin-ordinary}{n}}$ من المرات
  نحصل على !!{derivation} لـ~$!B$.
\end{proof}

سنشير إلى استعمال هذه القضية بـ~\PL.

\begin{prop}
  إذا كان $\Log{K} \Proves !A_{\OLMathNumeral{1}} \lif (!A_{\OLMathNumeral{2}} \lif \cdots (!A_{{\OLId{latin-ordinary}{n}}-{\OLMathNumeral{1}}} \lif
  !A_{\OLId{latin-ordinary}{n}})\dots)$، فعندئذ $\Log{K} \Proves \Box !A_{\OLMathNumeral{1}} \lif (\Box !A_{\OLMathNumeral{2}} \lif
  \cdots (\Box !A_{{\OLId{latin-ordinary}{n}}-{\OLMathNumeral{1}}} \lif \Box !A_{\OLId{latin-ordinary}{n}})\dots)$.
\end{prop}

\begin{proof}
  بالاستقراء على ${\OLId{latin-ordinary}{n}}$، تماما كما في برهان \olref[nor]{prop:rk}.
\end{proof}

سنشير إلى استعمال هذه القضية بـ~\RK. ولنوضح كيف تساعد هذه النتائج
على إثبات نتائج !!{derivability} بسهولة أكبر.

\begin{prop}
  $\Log{K} \Proves (\Box!A \land \Box!B) \lif \Box (!A \land !B)$
\end{prop}

\begin{proof}
  \begin{derivation}
    1. & $\Log{K} \Proves !A \lif (!B \lif (!A \land !B))$ & \Taut \\
    2. & $\Log{K} \Proves \Box!A \lif (\Box !B \lif \Box(!A \land !B))$ & \RK, 1\\
    3. & $\Log{K} \Proves (\Box!A \land \Box!B) \lif \Box(!A \land !B)$ & \PL, 2
  \end{derivation}
\end{proof}

% تصحيح مصدر (0432-I1): أُثبت واسم الصيغة في وسيط الإحلال؛ فقد سقطت علامة ! من B في الأصل الإنجليزي.
\begin{prop}\ollabel{prop:rewriting}
  إذا كان $\Log{K} \Proves !A \liff !B$ و$\Log{K} \Proves
  \Subst{!C}{!A}{{\OLId{latin-ordinary}{q}}}$، فعندئذ $\Log{K} \Proves \Subst{!C}{!B}{{\OLId{latin-ordinary}{q}}}$.
\end{prop}

\begin{proof}
  تمرين.
\end{proof}

\begin{prob}
  أثبت \olref[nml][prf][der]{prop:rewriting} بأن تثبت، بالاستقراء على
  تعقيد~$!C$، أنه إذا كان $\Log{K} \Proves !A \liff !B$، فعندئذ
  $\Log{K} \Proves \Subst{!C}{!A}{{\OLId{latin-ordinary}{q}}} \liff \Subst{!C}{!B}{{\OLId{latin-ordinary}{q}}}$.
\end{prob}

تفيد هذه القضية بوجه خاص عندما نريد تحويل $\Diamond$ إلى $\Box$
(أو العكس)، أو حذف النفي المزدوج من داخل !!a{formula}. وفيما يأتي،
سنعلّم تطبيقات \olref{prop:rewriting} بعبارة «$!B$ بدلا من $!A$»
كلما أعدنا كتابة !!a{formula}~$!C(!A)$ على صورة~$!C(!B)$. وبعبارة
أخرى، تختصر «$!B$ بدلا من $!A$» ما يأتي:
  \begin{derivation}
    & $\Proves !C(!A)$ \\
    & $\Proves !A \liff !B$\\
    & $\Proves !C(!B)$ & بواسطة \olref{prop:rewriting}
  \end{derivation}
فمثلا:

\begin{prop}
  $\Log{K} \Proves \lnot\Box {\OLId{latin-ordinary}{p}} \lif \Diamond \lnot {\OLId{latin-ordinary}{p}}$
\end{prop}

\begin{proof}
  \iftag{prvDiamond}{% Diamond is primitive
  \begin{derivation}
    1. & $\Log{K} \Proves \Diamond \lnot {\OLId{latin-ordinary}{p}} \liff \lnot\Box\lnot\lnot {\OLId{latin-ordinary}{p}}$ &
    \Dual\\
    2. & $\Log{K} \Proves \lnot \Box \lnot\lnot {\OLId{latin-ordinary}{p}} \lif \Diamond \lnot {\OLId{latin-ordinary}{p}}$ & \PL, 1\\
    3. & $\Log{K} \Proves \lnot\Box {\OLId{latin-ordinary}{p}} \lif \Diamond \lnot {\OLId{latin-ordinary}{p}}$ & ${\OLId{latin-ordinary}{p}}$ بدلا من $\lnot\lnot {\OLId{latin-ordinary}{p}}$\\
  \end{derivation}
  }{% Diamond is defined
  \begin{derivation}
    \OLClassicalDigits{1}. & $\Log{K} \Proves \lnot \Box \lnot\lnot {\OLId{latin-ordinary}{p}} \lif \Diamond \lnot {\OLId{latin-ordinary}{p}}$ & \Taut\\
    \OLClassicalDigits{2}. & $\Log{K} \Proves \lnot\Box {\OLId{latin-ordinary}{p}} \lif \Diamond \lnot {\OLId{latin-ordinary}{p}}$ & ${\OLId{latin-ordinary}{p}}$ بدلا من $\lnot\lnot {\OLId{latin-ordinary}{p}}$\\
  \end{derivation}
  الصيغة في السطر~${\OLMathNumeral{1}}$ مثال لـ$\lnot {\OLId{latin-ordinary}{q}} \lif \lnot {\OLId{latin-ordinary}{q}}$، بإحلال
  $\Box\lnot\lnot {\OLId{latin-ordinary}{p}}$ محل~${\OLId{latin-ordinary}{q}}$. والصيغتان $\Diamond\lnot {\OLId{latin-ordinary}{p}}$ و
  $\lnot\Box\lnot\lnot {\OLId{latin-ordinary}{p}}$ متطابقتان، لأن $\Diamond$ معرّف بأنه
  ~ $\lnot\Box\lnot$.}
\end{proof}

في !!{derivation} أعلاه، تختصر الخطوة الأخيرة «${\OLId{latin-ordinary}{p}}$ بدلا من
$\lnot\lnot {\OLId{latin-ordinary}{p}}$» ما يأتي:
  \begin{derivation}
    & $\Log{K} \Proves \lnot \Box \lnot\lnot {\OLId{latin-ordinary}{p}} \lif \Diamond \lnot
    {\OLId{latin-ordinary}{p}}$\\
    & $\Log{K} \Proves \lnot\lnot {\OLId{latin-ordinary}{p}} \liff {\OLId{latin-ordinary}{p}}$ & \Taut\\
    & $\Log{K} \Proves \lnot\Box {\OLId{latin-ordinary}{p}} \lif \Diamond \lnot {\OLId{latin-ordinary}{p}}$ & بواسطة \olref{prop:rewriting}
  \end{derivation}
وتقوم هنا الأدوار التي لـ$!C({\OLId{latin-ordinary}{q}})$ و$!A$ و$!B$ في
\olref{prop:rewriting}، على الترتيب، الصيغ $\lnot \Box {\OLId{latin-ordinary}{q}} \lif
\Diamond \lnot {\OLId{latin-ordinary}{p}}$ و$\lnot\lnot {\OLId{latin-ordinary}{p}}$ و~${\OLId{latin-ordinary}{p}}$.

عندما تحتوي !!a{formula} على !!{formula} جزئية $\lnot\Diamond !A$،
يمكننا استبدال $\Box\lnot !A$ بها باستعمال \olref{prop:rewriting}،
لأن $\Log{K} \Proves \lnot\Diamond !A \liff \Box\lnot !A$. وسنشير
إلى هذا الاستبدال وأمثاله ببساطة بعبارة «$\Box\lnot$ بدلا من
$\lnot\Diamond$».

ويبقى تسويغ الكلام على !!{derivability} بصيغة المخطط.
فالقضية السابقة لا تنحصر في~$\Log{K} \Proves \lnot\Box {\OLId{latin-ordinary}{p}} \lif
\Diamond \lnot {\OLId{latin-ordinary}{p}}$، بل تعطي
$\Log{K} \Proves \lnot\Box !A \lif \Diamond \lnot !A$
مهما كانت~$!A$. والقضية الآتية تبيّن وجه هذا التعميم.

\begin{prop}
  إذا كانت $!A$ مثال إحلال لـ$!B$ وكان $\Log{K} \Proves !B$، فعندئذ
  $\Log{K} \Proves !A$.
\end{prop}

\begin{proof}
  يُتحقق من ذلك بالاستقراء على طول !!{derivation} لـ~$!B$:
  إذا أُجري الإحلال في !!{derivation} كله بقي !!{derivation}
  صحيحًا. وتفصيله أن الإحلال في حالة تحصيل منطقي يعطي حالة
  تحصيل منطقي، والإحلال في حالة من
  \iftag{prvDiamond}{\Dual{} و~}{}\Ax{K} يعطي حالة من
  \iftag{prvDiamond}{\Dual{} و~}{}\Ax{K} أيضًا.
  وتبقى خطوتا~\MP{} و~\Nec{} صحيحتين إذا أُحلت !!{formula}s
  محل !!{propositional variable}s في المقدمات والنتيجة جميعًا.
  وليس في التحقق صعوبة جديدة، وإنما فيه تكرار الحالات.
\end{proof}


\end{document}
